\documentclass[11pt,a4paper]{article}

\usepackage[utf8]{inputenc}
\usepackage[T1]{fontenc}
\usepackage[top=3.5cm,bottom=3cm,left=2.5cm,right=2.5cm]{geometry}
\usepackage{tikz}
\usepackage[table]{xcolor}

\usetikzlibrary{positioning, arrows.meta, fit, backgrounds, calc, shadows}

\definecolor{PrimaryBlue}{RGB}{27,58,109}
\definecolor{PrimaryBlueLight}{RGB}{235,241,248}
\definecolor{AccentBlue}{RGB}{41,98,175}
\definecolor{DarkText}{RGB}{38,38,38}
\definecolor{LightGray}{RGB}{248,249,250}
\definecolor{MediumGray}{RGB}{118,118,118}
\definecolor{BorderGray}{RGB}{213,220,230}
\definecolor{SlateBlue}{RGB}{54,79,107}
\definecolor{SoftGreen}{RGB}{232,246,240}
\definecolor{SuccessGreen}{RGB}{34,139,84}
\definecolor{TableHead}{RGB}{240,245,252}
\definecolor{TableStripe}{RGB}{249,251,254}

\usepackage{amssymb}
\usepackage[scaled=0.92]{helvet}
\renewcommand{\familydefault}{\sfdefault}
\usepackage{microtype}
\usepackage{setspace}
\setstretch{1.15}
\usepackage{parskip}
\usepackage{graphicx}
\usepackage{tabularx}
\usepackage{booktabs}
\usepackage{array}
\usepackage{enumitem}
\usepackage{xurl}
\usepackage[hidelinks,breaklinks=true]{hyperref}
\usepackage{multirow}
\usepackage{ragged2e}
\usepackage{fontawesome5}

% ── FontAwesome5 icons ────────────────────────────────────────────────────────
% Standard (blue) — for use on white/light backgrounds
\newcommand{\iconResults}{\textcolor{AccentBlue}{\faChartBar}}
\newcommand{\iconServer}{\textcolor{AccentBlue}{\faServer}}
\newcommand{\iconInfo}{\textcolor{AccentBlue}{\faInfoCircle}}
\newcommand{\iconDir}{\textcolor{AccentBlue}{\faFolderOpen}}
\newcommand{\iconKey}{\textcolor{AccentBlue}{\faKey}}
\newcommand{\iconLock}{\textcolor{AccentBlue}{\faLock}}
\newcommand{\iconAlert}{\textcolor{AccentBlue}{\faExclamationTriangle}}
\newcommand{\iconTwoFA}{\textcolor{AccentBlue}{\faUserShield}}
\newcommand{\iconArrow}{\textcolor{AccentBlue}{\faArrowRight}}
\newcommand{\iconBug}{\textcolor{AccentBlue}{\faBug}}
\newcommand{\iconPerspective}{\textcolor{AccentBlue}{\faRocket}}
% White variants — for use on dark/blue backgrounds (e.g. PrimaryBlue title bars)
\newcommand{\iconResultsW}{\textcolor{white}{\faChartBar}}
\newcommand{\iconPerspectiveW}{\textcolor{white}{\faRocket}}
% ─────────────────────────────────────────────────────────────────────────────

\usepackage[skins]{tcolorbox}
\tcbset{before skip=0.6em,after skip=0.6em}

\usepackage{titlesec}
\titleformat{\section}
  {\color{PrimaryBlue}\Large\bfseries}
  {}
  {0pt}
  {}[\vspace{-0.3ex}\textcolor{BorderGray}{\rule{\textwidth}{1pt}}]

\titleformat{\subsection}
  {\color{SlateBlue}\normalsize\bfseries}
  {}
  {0pt}
  {}

\titlespacing{\section}{0pt}{1.4em}{0.5em}
\titlespacing{\subsection}{0pt}{0.9em}{0.3em}

\usepackage{fancyhdr}
\pagestyle{fancy}
\fancyhf{}
\renewcommand{\headrulewidth}{0pt}
\setlength{\headheight}{38pt}
\setlength{\headsep}{12pt}
\lhead{%
  \begin{tikzpicture}[baseline=(T.base)]
    \node[inner sep=0pt] (T) {\textcolor{MediumGray}{\small\textbf{DOCUMENTATION TECHNIQUE FINALE}}};
  \end{tikzpicture}%
}
\rhead{\textcolor{PrimaryBlue}{\small\bfseries Guardia School}}
\lfoot{\textcolor{MediumGray}{\footnotesize Document Confidentiel \textperiodcentered{} Usage Interne}}
\rfoot{\textcolor{PrimaryBlue}{\small\textbf{\thepage}}}

\hypersetup{
  colorlinks=true,
  linkcolor=DarkText,
  urlcolor=AccentBlue,
  pdftitle={Guardia School - Documentation Technique Finale}
}

\setlist[itemize]{leftmargin=1.4em,itemsep=0.2em,topsep=0.2em,
  label={\textcolor{AccentBlue}{\small\textbullet}}}
\setlength{\parskip}{0.5em}

\newcommand{\statusdone}{%
  \textcolor{SuccessGreen}{\small\faCheck\hspace{0.3em}\textbf{Terminé}}%
}

\newcolumntype{L}[1]{>{\raggedright\arraybackslash}p{#1}}
\newcolumntype{C}[1]{>{\centering\arraybackslash}p{#1}}




\begin{document}
\sloppy

\hypersetup{pageanchor=false}
\begin{titlepage}
  \thispagestyle{empty}
  \begin{tikzpicture}[remember picture,overlay]
    \fill[PrimaryBlue] (current page.north west) rectangle ++(1.6cm,-\paperheight);
    \fill[PrimaryBlueLight] (current page.north east) ++(0,-\paperheight) rectangle ++(-6cm,6cm);
  \end{tikzpicture}

  \vspace*{1.2cm}
  \hspace*{1.2cm}
  \begin{minipage}{0.8\textwidth}
    \textcolor{PrimaryBlue}{\fontsize{28}{34}\selectfont\bfseries Guardia School}
  \end{minipage}

  \vspace{3.5cm}

  \hspace*{1.2cm}{\fontsize{22}{28}\selectfont\bfseries\textcolor{DarkText}{Documentation Technique Finale}}\par\vspace{0.6cm}
  \hspace*{1.2cm}{\large\textcolor{AccentBlue}{Déploiement Automatisé d'une Infrastructure Virtualisée Sécurisée}}\par\vspace{1.4cm}

  \hspace*{1.2cm}\textcolor{PrimaryBlue}{\rule{6cm}{2.5pt}}\par\vspace{1.4cm}

  \hspace*{1.2cm}{\normalsize\textcolor{MediumGray}{Ansible \textperiodcentered{} SNMP \textperiodcentered{} HashiCorp Vault \textperiodcentered{} UFW \textperiodcentered{} 2FA \textperiodcentered{} Supervision automatique}}\par\vspace{2.8cm}

  \hspace*{1.2cm}
  \begin{minipage}{0.75\textwidth}
    \begin{tcolorbox}[
      enhanced,
      colback=white,
      colframe=BorderGray,
      boxrule=0.8pt,
      arc=8pt,
      left=14pt,right=14pt,top=12pt,bottom=12pt
    ]
    \renewcommand{\arraystretch}{1.5}
    \begin{tabular}{@{} L{3.8cm} L{6cm} @{}}
      \textcolor{MediumGray}{\small Projet} & \textbf{Infrastructure Ansible Sécurisée} \\
      \textcolor{MediumGray}{\small Client} & \textbf{Guardia School} \\
      \textcolor{MediumGray}{\small Équipe} & \textbf{Équipe Projet Ansible} \\
      \textcolor{MediumGray}{\small Référence} & \texttt{\small GS-ANSIBLE-SEC-2026-F} \\
      \textcolor{MediumGray}{\small Période} & \textbf{Avril 2026} \\
      \textcolor{MediumGray}{\small Version} & \textbf{1.0 Final} \\
      \textcolor{MediumGray}{\small Finalisé le} & \textbf{23 avril 2026} \\
    \end{tabular}
    \end{tcolorbox}
  \end{minipage}

  \vfill
  \hspace*{1.2cm}{\small\textit{\textcolor{MediumGray}{Document Confidentiel \textperiodcentered{} Usage Interne}}}
\end{titlepage}
\hypersetup{pageanchor=true}

\clearpage
\renewcommand{\contentsname}{\textcolor{PrimaryBlue}{Sommaire}}
{\setstretch{1.35}
\tableofcontents
}

\clearpage

\section{Résumé Exécutif}

Ce document constitue la documentation technique finale du projet d'automatisation et de déploiement d'une infrastructure virtualisée sécurisée, réalisé par l'équipe IT de Guardia School.

L'objectif principal était de concevoir, déployer et documenter un système entièrement automatisé permettant de configurer des machines Linux  \textit{from scratch}, en intégrant les bonnes pratiques de sécurité : gestion centralisée des secrets, surveillance temps réel par SNMP, pare-feu UFW, authentification à deux facteurs (2FA), durcissement OS et inventaire automatique.

\begin{tcolorbox}[
  enhanced,
  colback=PrimaryBlueLight,
  colframe=PrimaryBlue,
  boxrule=1.2pt,
  arc=8pt,
  left=14pt,right=14pt,top=10pt,bottom=10pt,
  title={\iconResultsW\hspace{0.4em}\textbf{Résultats obtenus}},
  fonttitle=\bfseries\normalsize,
  coltitle=white,
  attach boxed title to top left={yshift=-2mm,xshift=6mm},
  boxed title style={colback=PrimaryBlue,arc=4pt,boxrule=0pt}
]
\begin{itemize}
  \item Déploiement complet en moins de \textbf{30 minutes} sur une machine vierge via playbooks Ansible
  \item \textbf{4 VMs} déployées et configurées (Linux)
  \item SNMP opérationnel sur toutes les cibles, intégré à Centreon
  \item HashiCorp Vault déployé avec rotation automatique des secrets
  \item UFW configuré avec politique \texttt{deny-all} et règles explicites
  \item 2FA (TOTP / Google Authenticator) activé sur tous les accès SSH critiques
  \item Inventaire automatique généré via Ansible
  \item Supervision automatique Centreon avec alertes e-mail
\end{itemize}
\end{tcolorbox}

\section{Introduction et Objectifs}

\subsection{Contexte du projet}

Dans un contexte de multiplication des cybermenaces et de complexité croissante des infrastructures IT, la gestion manuelle des configurations devient une source majeure de risques : erreurs humaines, inconsistances entre machines, délais de déploiement et difficultés de traçabilité.

Ce projet répond à ces enjeux par l'automatisation systématique, en s'appuyant sur Ansible, un outil d'automatisation IT agentless, open source, largement adopté en production.

\subsection{Objectifs techniques}
\begin{itemize}
  \item Déployer une infrastructure virtualisée complète sur Proxmox V
  \item Automatiser la configuration initiale (hardening) de toutes les machines via Ansible
  \item Mettre en place SNMP pour la remontée de métriques vers un serveur de supervision Centreon
  \item Déployer HashiCorp Vault comme gestionnaire centralisé de secrets
  \item Configurer UFW avec une politique restrictive de type \texttt{deny-all}
  \item Implémenter le 2FA sur les accès SSH
  \item Automatiser la rotation et le changement des mots de passe SSH
  \item Générer un inventaire automatique de toutes les machines
  \item Configurer la supervision automatique avec alertes en temps réel
\end{itemize}

\subsection{Périmètre technique}

\begin{tcolorbox}[
  enhanced,
  colback=LightGray,
  colframe=BorderGray,
  boxrule=0.8pt,
  arc=6pt,
  left=10pt,right=10pt,top=8pt,bottom=8pt
]
\renewcommand{\arraystretch}{1.4}
\begin{tabular}{@{} L{4.2cm} L{9cm} @{}}
  \textcolor{PrimaryBlue}{\textbf{Hyperviseur}} & Proxmox VE \\
  \textcolor{PrimaryBlue}{\textbf{Automatisation}} & Ansible  \\
  \textcolor{PrimaryBlue}{\textbf{Cibles Linux}} & Ubuntu  \\
  \textcolor{PrimaryBlue}{\textbf{Pare-feu / Routeur}} & OPNSense \\
  \textcolor{PrimaryBlue}{\textbf{Supervision}} & Centreon  \\
  \textcolor{PrimaryBlue}{\textbf{Gestion des secrets}} & HashiCorp Vault  \\
  \textcolor{PrimaryBlue}{\textbf{2FA}} & Google Authenticator (libpam-google-authenticator) \\
\end{tabular}
\end{tcolorbox}

\section{Infrastructure Virtualisée}

\subsection{Vue d'ensemble}

L'infrastructure repose sur un hyperviseur Proxmox hébergeant 4 machines virtuelles. Chaque VM remplit un rôle précis dans l'architecture. Le réseau est segmenté en deux zones : le réseau d'administration (\texttt{10.1.91.0/24}) et un réseau DMZ optionnel.

\begin{tcolorbox}[
  enhanced,
  colback=white,
  colframe=BorderGray,
  boxrule=0.8pt,
  arc=8pt,
  left=0pt,right=0pt,top=0pt,bottom=0pt,
  title={\iconServer\hspace{0.4em}\textbf{Tableau détaillé des VMs}},
  fonttitle=\bfseries\small,
  coltitle=PrimaryBlue,
  attach boxed title to top left={yshift=-2mm,xshift=6mm},
  boxed title style={colback=white,colframe=BorderGray,arc=4pt,boxrule=0.8pt}
]
\renewcommand{\arraystretch}{1.4}
\begin{tabularx}{\linewidth}{@{\hspace{8pt}} L{3.0cm} X C{0.9cm} C{1.1cm} C{2.1cm} L{3.0cm} @{\hspace{8pt}}}
  \rowcolor{TableHead}
  \textbf{\textcolor{PrimaryBlue}{VM / Hôte}} &
  \textbf{\textcolor{PrimaryBlue}{Rôle}} &
  \textbf{\textcolor{PrimaryBlue}{CPU}} &
  \textbf{\textcolor{PrimaryBlue}{RAM}} &
  \textbf{\textcolor{PrimaryBlue}{Stockage}} &
  \textbf{\textcolor{PrimaryBlue}{OS}} \\
  \midrule
  \texttt{srv-ansible} & Contrôleur Ansible & 2 & 2 Go & 20 Go SSD & Ubuntu 22.04 \\
  \rowcolor{TableStripe}
  \texttt{srv-monitoring} & Centreon / SNMP & 4 & 4 Go & 50 Go SSD & Ubuntu 22.04 \\
  \rowcolor{TableStripe}
  \texttt{linux-target-01} & Cible Linux & 2 & 2 Go & 30 Go SSD & Debian 12 \\
  \texttt{OPNSense-fw} & Pare-feu / Routeur & 2 & 1 Go & 16 Go SSD & OPNSense  \\
  \bottomrule
\end{tabularx}
\vspace{4pt}
\end{tcolorbox}

\subsection{Configuration réseau}

Toutes les VMs sont connectées au bridge \texttt{vmbr0} de Proxmox (réseau \texttt{10.1.91.0/24}). OPNSense fait office de routeur et de pare-feu périmétrique. Les IPs sont statiques, configurées dans les fichiers \texttt{/etc/netplan/} (Linux) ou via les propriétés de la carte réseau (Windows).

\subsection{ISOs utilisées}
\begin{itemize}
  \item Ubuntu 22.04.3 LTS Server \quad \texttt{ubuntu-22.04.3-live-server-amd64.iso} \hfill{\small 2.1 Go}
  \item Debian 12.4 Netinstall \quad\quad\quad \texttt{debian-12.4.0-amd64-netinst.iso} \hfill{\small 659 Mo}
  \item OPNSense  2.6.0 \quad\quad\quad\quad \texttt{OPNSense-2.6.0-RELEASE-amd64.iso} \hfill{\small 1.0 Go}
\end{itemize}

\section{Méthodologie et Architecture Ansible}

\subsection{Principe de fonctionnement}

Ansible est un outil d'automatisation agentless (sans agent à installer sur les cibles). Le contrôleur (\texttt{srv-ansible}) se connecte aux machines cibles via SSH (Linux) ou  et exécute des tâches définies dans des fichiers YAML appelés \textit{playbooks}.

\begin{tcolorbox}[
  enhanced,
  colback=PrimaryBlueLight,
  colframe=AccentBlue,
  boxrule=0pt,
  borderline west={3pt}{0pt}{AccentBlue},
  arc=0pt,
  left=10pt,right=10pt,top=7pt,bottom=7pt
]
\iconInfo\hspace{0.4em}\textbf{Avantage clé :} aucun logiciel n'est nécessaire sur les machines cibles Linux. Seuls SSH et Python\,3 sont requis.
\end{tcolorbox}

\subsection{Structure du projet Ansible}

\begin{tcolorbox}[
  enhanced,
  colback=LightGray,
  colframe=BorderGray,
  boxrule=0.8pt,
  arc=8pt,
  fontupper=\ttfamily\small,
  left=12pt,right=12pt,top=10pt,bottom=10pt,
  title={\iconDir\hspace{0.4em}\textbf{Arborescence du projet}},
  fonttitle=\bfseries\small,
  coltitle=PrimaryBlue,
  attach boxed title to top left={yshift=-2mm,xshift=6mm},
  boxed title style={colback=LightGray,colframe=BorderGray,arc=4pt,boxrule=0.8pt}
]
ansible-project/\\
\textcolor{PrimaryBlue}{|} ansible.cfg \hfill \normalfont\small\textcolor{MediumGray}{Configuration globale Ansible}\\
\textcolor{PrimaryBlue}{|} inventory/\\
\textcolor{PrimaryBlue}{|}\hspace{1em}\textcolor{PrimaryBlue}{|} hosts.yml \hfill \normalfont\small\textcolor{MediumGray}{Inventaire statique des hôtes}\\
\textcolor{PrimaryBlue}{|}\hspace{1em}\textcolor{PrimaryBlue}{+} group\_vars/\\
\textcolor{PrimaryBlue}{|}\hspace{2em}\textcolor{PrimaryBlue}{|} all.yml \hfill \normalfont\small\textcolor{MediumGray}{Variables globales}\\
\textcolor{PrimaryBlue}{|}\hspace{2em}\textcolor{PrimaryBlue}{|} linux.yml \hfill \normalfont\small\textcolor{MediumGray}{Variables Linux}\\

\textcolor{PrimaryBlue}{|} roles/\\
\textcolor{PrimaryBlue}{|}\hspace{1em}\textcolor{PrimaryBlue}{|} common/ \hfill \normalfont\small\textcolor{MediumGray}{Tâches communes (users, updates)}\\
\textcolor{PrimaryBlue}{|}\hspace{1em}\textcolor{PrimaryBlue}{|} hardening/ \hfill \normalfont\small\textcolor{MediumGray}{Durcissement SSH, OS}\\
\textcolor{PrimaryBlue}{|}\hspace{1em}\textcolor{PrimaryBlue}{|} ufw/ \hfill \normalfont\small\textcolor{MediumGray}{Configuration pare-feu UFW}\\
\textcolor{PrimaryBlue}{|}\hspace{1em}\textcolor{PrimaryBlue}{|} snmp/ \hfill \normalfont\small\textcolor{MediumGray}{Déploiement SNMP agent}\\
\textcolor{PrimaryBlue}{|}\hspace{1em}\textcolor{PrimaryBlue}{|} vault\_client/ \hfill \normalfont\small\textcolor{MediumGray}{Intégration HashiCorp Vault}\\
\textcolor{PrimaryBlue}{|}\hspace{1em}\textcolor{PrimaryBlue}{|} 2fa/ \hfill \normalfont\small\textcolor{MediumGray}{Configuration 2FA SSH}\\
\textcolor{PrimaryBlue}{|}\hspace{1em}\textcolor{PrimaryBlue}{|} Centreon\_agent/ \hfill \normalfont\small\textcolor{MediumGray}{Agent Centreon}\\
\textcolor{PrimaryBlue}{|} playbooks/\\
\textcolor{PrimaryBlue}{|}\hspace{1em}\textcolor{PrimaryBlue}{|} site.yml \hfill \normalfont\small\textcolor{MediumGray}{Playbook principal (full deploy)}\\
\textcolor{PrimaryBlue}{|}\hspace{1em}\textcolor{PrimaryBlue}{|} hardening.yml \hfill \normalfont\small\textcolor{MediumGray}{Hardening seul}\\
\textcolor{PrimaryBlue}{|}\hspace{1em}\textcolor{PrimaryBlue}{|} inventory\_report.yml \hfill \normalfont\small\textcolor{MediumGray}{Génération inventaire}\\
\textcolor{PrimaryBlue}{|}\hspace{1em}\textcolor{PrimaryBlue}{+} rotate\_passwords.yml \hfill \normalfont\small\textcolor{MediumGray}{Rotation mots de passe SSH}\\
\textcolor{PrimaryBlue}{+} vault/\\
\hspace{1em}\textcolor{PrimaryBlue}{+} secrets.yml \hfill \normalfont\small\textcolor{MediumGray}{Secrets chiffrés (ansible-vault)}
\end{tcolorbox}

\section{Inventaire Automatique}

\subsection{Inventaire statique (hosts.yml)}

Le fichier d'inventaire référence toutes les machines avec leurs paramètres de connexion SSH (IP, port, utilisateur, clé). Il constitue la base de données de l'infrastructure.

\subsection{Tableau de l'inventaire}

\begin{tcolorbox}[
  enhanced,
  colback=white,
  colframe=BorderGray,
  boxrule=0.8pt,
  arc=8pt,
  left=0pt,right=0pt,top=0pt,bottom=0pt
]
\renewcommand{\arraystretch}{1.4}
\begin{tabularx}{\linewidth}{@{\hspace{8pt}} L{3.3cm} L{2.4cm} C{1.0cm} C{1.5cm} X @{\hspace{8pt}}}
  \rowcolor{TableHead}
  \textbf{\textcolor{PrimaryBlue}{Hôte}} &
  \textbf{\textcolor{PrimaryBlue}{IP}} &
  \textbf{\textcolor{PrimaryBlue}{Port}} &
  \textbf{\textcolor{PrimaryBlue}{User}} &
  \textbf{\textcolor{PrimaryBlue}{Rôles attribués}} \\
  \midrule
  \texttt{srv-ansible} & \texttt{10.1.91.100} & 22 & ansible & controller \\
  \rowcolor{TableStripe}
  \texttt{srv-monitoring} & \texttt{10.1.91.104} & 22 & ansible & Centreon, snmp, hardening \\
  \texttt{linux-target-01} & \texttt{10.1.91.102} & 22 & ansible & common, ufw, snmp-agent, hardening \\
  \rowcolor{TableStripe}
  \texttt{OPNSense} & \texttt{10.1.91.114} & 22 & ansible & gateaway \\
  \bottomrule
\end{tabularx}
\vspace{4pt}
\end{tcolorbox}

\subsection{Inventaire automatique via playbook}

Le module \texttt{setup} d'Ansible collecte automatiquement les facts de chaque machine : IP, OS, version du kernel, RAM, CPU, disques, interfaces réseau, etc. Le playbook associé génère un rapport complet en JSON et HTML.

\section{Hardening SSH et Gestion des Accès}

\subsection{Automatisation du hardening SSH}

Le rôle \texttt{hardening} applique une configuration sécurisée du démon SSH sur toutes les machines Linux. Les principales actions sont : changement de port, désactivation du login root, restriction aux clés SSH et affichage d'une bannière d'avertissement légal.

\subsection{Rotation automatique des mots de passe SSH}

Un playbook dédié automatise le changement de mot de passe des utilisateurs sur toutes les machines, en récupérant les nouveaux secrets depuis HashiCorp Vault.

\section{Authentification à Deux Facteurs (2FA)}

\subsection{Présentation}

Le 2FA est implémenté via le module PAM Google Authenticator (\texttt{libpam-google-authenticator}). Chaque connexion SSH nécessite deux éléments combinés : la clé SSH (ce que l'on possède) et un code TOTP à 6 chiffres généré toutes les 30 secondes (ce que l'on sait).



\subsection{Déploiement via Ansible}

Après le déploiement, chaque administrateur doit exécuter \texttt{google-authenticator} sur son compte pour générer le QR code à scanner dans Google Authenticator ou FreeOTP.

\section{Pare-feu UFW}

\subsection{Politique de sécurité}

UFW (Uncomplicated Firewall) est le front-end iptables utilisé sur toutes les machines Linux. La politique appliquée est \texttt{deny-all} par défaut : tout trafic est bloqué sauf ce qui est explicitement autorisé. Cette approche de liste blanche minimise la surface d'attaque.

\section{Supervision SNMP}

\subsection{Architecture SNMP}

SNMP (Simple Network Management Protocol) est déployé en mode agent sur chaque machine cible. Le serveur Centreon interroge ces agents toutes les 60 secondes pour collecter les métriques.

\begin{tcolorbox}[
  enhanced,
  colback=PrimaryBlueLight,
  colframe=AccentBlue,
  boxrule=0pt,
  borderline west={3pt}{0pt}{AccentBlue},
  arc=0pt,
  left=10pt,right=10pt,top=7pt,bottom=7pt
]
\iconLock\hspace{0.4em}\textbf{Sécurité SNMP :} authentification SHA-256 et chiffrement AES-256. Chaque agent dispose de credentials uniques stockés dans HashiCorp Vault.
\end{tcolorbox}



\section{HashiCorp Vault}

\subsection{Rôle de Vault dans l'architecture}

HashiCorp Vault est le coffre-fort centralisé de l'infrastructure. Il stocke de façon sécurisée tous les secrets : mots de passe SSH, credentials SNMP, tokens API, certificats TLS. Ansible récupère ces secrets dynamiquement à l'exécution, évitant de stocker des informations sensibles en clair dans les fichiers de configuration.

\begin{tcolorbox}[
  enhanced,
  colback=white,
  colframe=BorderGray,
  boxrule=0.8pt,
  arc=8pt,
  left=10pt,right=10pt,top=10pt,bottom=10pt,
  title={\iconKey\hspace{0.4em}\textbf{Caractéristiques de sécurité}},
  fonttitle=\bfseries\small,
  coltitle=PrimaryBlue,
  attach boxed title to top left={yshift=-2mm,xshift=6mm},
  boxed title style={colback=white,colframe=BorderGray,arc=4pt,boxrule=0.8pt}
]
\begin{itemize}
  \item Chiffrement \textbf{AES-256-GCM} des données au repos
  \item Authentification par token, AppRole ou certificat TLS
  \item Audit log complet de tous les accès
  \item Rotation automatique des secrets (lease TTL)
  \item Interface web HTTPS, API REST et CLI
\end{itemize}
\end{tcolorbox}

\section{Supervision Automatique}

\subsection{Architecture de supervision}

La supervision repose sur Centreon installé sur \texttt{srv-monitoring}. Centreon interroge les agents Centreon et SNMP sur toutes les machines cibles. Des seuils d'alerte sont configurés pour envoyer des notifications par e-mail en cas de dépassement.

\begin{tcolorbox}[
  enhanced,
  colback=LightGray,
  colframe=BorderGray,
  boxrule=0.8pt,
  arc=6pt,
  left=10pt,right=10pt,top=8pt,bottom=8pt
]
\renewcommand{\arraystretch}{1.35}
\begin{tabular}{@{} L{4.5cm} L{9cm} @{}}
  \textcolor{PrimaryBlue}{\textbf{Métriques collectées}} & CPU, RAM, disque, réseau, services actifs, uptime \\
  \textcolor{PrimaryBlue}{\textbf{Fréquence de polling}} & 60 s (métriques normales) / 5 s (alertes critiques) \\
  \textcolor{PrimaryBlue}{\textbf{Rétention des données}} & 30 jours en base MariaDB \\
  \textcolor{PrimaryBlue}{\textbf{Seuils d'alerte}} & CPU > 85\%, RAM > 90\%, disque > 80\% \\
\end{tabular}
\end{tcolorbox}

\subsection{Auto-supervision Centreon}

Centreon supervise également son propre serveur (\texttt{srv-monitoring}) grâce au template \textit{Linux by Centreon agent} appliqué en boucle locale. Les métriques du serveur de supervision sont ainsi visibles dans Centreon lui-même, permettant de détecter toute défaillance du service de supervision.

\section{Déploiement sur Machine Vierge}

\subsection{Prérequis}

Cette procédure permet de déployer l'intégralité de l'infrastructure à partir d'une machine Ubuntu 22.04 fraîchement installée, en moins de 30 minutes.

\begin{itemize}
  \item Machine avec Debian installé 
  \item Accès réseau fonctionnel (\texttt{10.1.91.0/24})
  \item Accès SSH root ou sudo vers les cibles
  \item Git installé : \texttt{sudo apt install git -y}
\end{itemize}

\subsection{Procédure complète de déploiement}

\begin{tcolorbox}[
  enhanced,
  colback=white,
  colframe=BorderGray,
  boxrule=0.8pt,
  arc=8pt,
  left=10pt,right=10pt,top=10pt,bottom=10pt
]
\renewcommand{\arraystretch}{1.5}
\begin{tabular}{@{} C{1.2cm} L{13cm} @{}}
  \textcolor{white}{\colorbox{PrimaryBlue}{\textbf{\small\ 1\ }}} & Cloner le projet Ansible \\
  \textcolor{white}{\colorbox{PrimaryBlue}{\textbf{\small\ 2\ }}} & Lancer le script d'initialisation du contrôleur \\
  \textcolor{white}{\colorbox{PrimaryBlue}{\textbf{\small\ 3\ }}} & Distribuer la clé SSH sur les cibles \\
  \textcolor{white}{\colorbox{PrimaryBlue}{\textbf{\small\ 4\ }}} & Vérifier la connectivité avec toutes les machines \\
  \textcolor{white}{\colorbox{PrimaryBlue}{\textbf{\small\ 5\ }}} & Déploiement complet via le playbook principal \\
  \textcolor{white}{\colorbox{PrimaryBlue}{\textbf{\small\ 6\ }}} & Génération de l'inventaire automatique \\
  \textcolor{white}{\colorbox{PrimaryBlue}{\textbf{\small\ 7\ }}} & Vérification finale \\
\end{tabular}
\end{tcolorbox}

\subsection{Temps de déploiement mesuré}

\begin{tcolorbox}[
  enhanced,
  colback=white,
  colframe=BorderGray,
  boxrule=0.8pt,
  arc=8pt,
  left=0pt,right=0pt,top=0pt,bottom=0pt
]
\renewcommand{\arraystretch}{1.4}
\begin{tabularx}{\linewidth}{@{\hspace{8pt}} X C{2.8cm} L{3.5cm} @{\hspace{8pt}}}
  \rowcolor{TableHead}
  \textbf{\textcolor{PrimaryBlue}{Phase}} &
  \textbf{\textcolor{PrimaryBlue}{Durée moyenne}} &
  \textbf{\textcolor{PrimaryBlue}{Machines}} \\
  \midrule
  Bootstrap contrôleur & 3 min & 1 (srv-ansible) \\
  \rowcolor{TableStripe}
  Distribution clés SSH & 2 min & 3 cibles \\
  Hardening + UFW & 8 min & 3 cibles \\
  \rowcolor{TableStripe}
  SNMP + Centreon Agent & 6 min & 3 cibles \\
  Vault + 2FA & 5 min & 2 machines critiques \\
  \rowcolor{TableStripe}
  Inventaire + Vérification & 4 min & Toutes \\
  \midrule
  \rowcolor{PrimaryBlueLight}
  \textbf{TOTAL} & \textbf{{\raise.15ex\hbox{$\sim$}}28 min} & \textbf{5 VMs} \\
  \bottomrule
\end{tabularx}
\vspace{4pt}
\end{tcolorbox}

\section{Avancement Chronologique du Projet}

Cette section documente l'avancement jour par jour, permettant de comprendre les choix techniques effectués, les obstacles rencontrés et les solutions apportées.

\begin{tcolorbox}[
  enhanced,
  colback=white,
  colframe=BorderGray,
  boxrule=0.8pt,
  arc=8pt,
  left=0pt,right=0pt,top=0pt,bottom=0pt
]
\renewcommand{\arraystretch}{1.5}
\begin{tabularx}{\linewidth}{@{\hspace{8pt}} C{0.8cm} C{2.4cm} X C{2.2cm} @{\hspace{8pt}}}
  \rowcolor{TableHead}
  \textbf{\textcolor{PrimaryBlue}{Jour}} &
  \textbf{\textcolor{PrimaryBlue}{Date}} &
  \textbf{\textcolor{PrimaryBlue}{Actions réalisées}} &
  \textbf{\textcolor{PrimaryBlue}{Statut}} \\
  \midrule
  \textbf{J1} & 20 Avr. 2026 &
    Initialisation Proxmox, création des VMs, configuration réseau, attribution des IPs statiques, installation Ubuntu 22.04 et Debian 12
    & \statusdone \\
  \rowcolor{TableStripe}
  \textbf{J2} & 21 Avr. 2026 &
    Installation Ansible sur \texttt{srv-ansible}, génération des clés SSH, déploiement de l'inventaire initial, tests de connectivité
    & \statusdone \\
  \textbf{J3} & 22 Avr. 2026 &
    Playbooks hardening SSH (désactivation root, changement de port, bannière), déploiement UFW, configuration SNMP , intégration Centreon
    & \statusdone \\
  \rowcolor{TableStripe}
  \textbf{J4} & 23 Avr. 2026 &
    Déploiement HashiCorp Vault , intégration Vault + Ansible, mise en place 2FA (Google Authenticator via PAM), rotation automatique des mots de passe SSH
    & \statusdone \\
  \textbf{J5} & 24 Avr. 2026 &
    Inventaire automatique via module \texttt{setup}, génération rapports ,  rédaction du rapport final, revue de sécurité, démonstration du déploiement complet
    & \statusdone \\
  \bottomrule
\end{tabularx}
\vspace{4pt}
\end{tcolorbox}

\subsection{Obstacles rencontrés et solutions}

\begin{tcolorbox}[
  enhanced,
  colback=white,
  colframe=BorderGray,
  boxrule=0.8pt,
  arc=8pt,
  borderline west={4pt}{0pt}{PrimaryBlue},
  left=10pt,right=10pt,top=8pt,bottom=8pt,
  title={\iconBug\hspace{0.3em}SNMP : credentials non synchronisés après rotation Vault},
  fonttitle=\bfseries\small,
  coltitle=PrimaryBlue,
  attach boxed title to top left={yshift=-2mm,xshift=6mm},
  boxed title style={colback=white,colframe=BorderGray,arc=4pt,boxrule=0.8pt}
]
Lors des premières rotations de credentials SNMP via Vault, les fichiers \texttt{snmpd.conf} n'étaient pas automatiquement mis à jour. La solution : ajouter un handler Ansible déclenché uniquement en cas de changement, couplé à un re-déploiement via \texttt{rotate\_passwords.yml}.
\end{tcolorbox}

\begin{tcolorbox}[
  enhanced,
  colback=white,
  colframe=BorderGray,
  boxrule=0.8pt,
  arc=8pt,
  borderline west={4pt}{0pt}{PrimaryBlue},
  left=10pt,right=10pt,top=8pt,bottom=8pt,
  title={\iconBug\hspace{0.3em}Le 2FA bloquait les connexions Ansible},
  fonttitle=\bfseries\small,
  coltitle=PrimaryBlue,
  attach boxed title to top left={yshift=-2mm,xshift=6mm},
  boxed title style={colback=white,colframe=BorderGray,arc=4pt,boxrule=0.8pt}
]
L'activation du 2FA sur \texttt{srv-ansible} a initialement bloqué les connexions Ansible. La solution : créer un compte de service \texttt{ansible-svc} avec authentification par clé SSH uniquement (\texttt{AuthenticationMethods publickey}) et exclure ce compte du 2FA PAM via une directive \texttt{Match User} dans \texttt{sshd\_config}.
\end{tcolorbox}

\section{État Final et Recommandations}

\subsection{État final de l'infrastructure}

À l'issue du projet, l'ensemble des objectifs définis ont été atteints :

\begin{tcolorbox}[
  enhanced,
  colback=white,
  colframe=BorderGray,
  boxrule=0.8pt,
  arc=8pt,
  left=0pt,right=0pt,top=0pt,bottom=0pt
]
\renewcommand{\arraystretch}{1.4}
\begin{tabularx}{\linewidth}{@{\hspace{8pt}} L{3.8cm} X C{2.2cm} @{\hspace{8pt}}}
  \rowcolor{TableHead}
  \textbf{\textcolor{PrimaryBlue}{Composant}} &
  \textbf{\textcolor{PrimaryBlue}{État final}} &
  \textbf{\textcolor{PrimaryBlue}{Statut}} \\
  \midrule
  Ansible & Opérationnel, clés SSH distribuées, inventaire actif & \statusdone \\
  \rowcolor{TableStripe}
  UFW Firewall & Deny-all actif, règles SSH / SNMP / Centreon configurées & \statusdone \\
  SNMP & Agents opérationnels sur 3 cibles Linux, auth SHA + AES & \statusdone \\
  \rowcolor{TableStripe}
  HashiCorp Vault & Unsealed, KV v2 actif, rotation automatique configurée & \statusdone \\
  Supervision Centreon & 3 hôtes supervisés& \statusdone \\
  \rowcolor{TableStripe}
  Hardening SSH & Port 2222, root désactivé, clés seules, bannière active & \statusdone \\
  2FA SSH & TOTP Google Auth actif sur les 3 machines critiques & \statusdone \\
  \rowcolor{TableStripe}
  Inventaire automatique & (IP, port, OS, CPU, RAM) & \statusdone \\
  \bottomrule
\end{tabularx}
\vspace{4pt}
\end{tcolorbox}

\subsection{Recommandations post-projet}

\begin{tcolorbox}[
  enhanced,
  colback=white,
  colframe=BorderGray,
  boxrule=0.8pt,
  arc=8pt,
  left=10pt,right=10pt,top=10pt,bottom=10pt,
  title={\iconArrow\hspace{0.4em}\textbf{Actions recommandées}},
  fonttitle=\bfseries\small,
  coltitle=PrimaryBlue,
  attach boxed title to top left={yshift=-2mm,xshift=6mm},
  boxed title style={colback=white,colframe=BorderGray,arc=4pt,boxrule=0.8pt}
]
\begin{itemize}
\RaggedRight
  \item \textbf{Vault en production :} activer l'auto-unseal via un HSM ou AWS KMS / Azure Key Vault pour éviter l'unseal manuel à chaque redémarrage.
  \item \textbf{Centralisation des logs :} déployer un SIEM (Elastic Stack, Wazuh) pour agréger les logs Ansible, Vault audit, SSH et UFW.
  \item \textbf{Rotation régulière des secrets :} programmer une tâche cron pour exécuter \texttt{rotate\_passwords.yml} automatiquement chaque semaine.
  \item \textbf{Tests de restauration :} valider mensuellement la capacité à redéployer l'infrastructure \textit{from scratch} sur une VM vierge.
  \item \textbf{Mise à jour Ansible :} vérifier mensuellement les nouvelles versions d'Ansible et des collections (\texttt{ansible-galaxy collection list}).
\end{itemize}
\end{tcolorbox}

\section{Conclusion}


Ce projet a démontré la puissance de l'automatisation via Ansible pour déployer et sécuriser une infrastructure multi-systèmes (Linux ) de façon cohérente, reproductible et rapide.

En moins de 30 minutes, il est désormais possible de configurer intégralement une machine vierge avec : hardening SSH, pare-feu UFW, agents SNMP, supervision Centreon, intégration HashiCorp Vault et authentification 2FA. Ce qui nécessiterait plusieurs heures de travail manuel peut être accompli en quelques commandes, sans erreur humaine.

La documentation présente dans ce rapport constitue une référence complète : toute personne disposant des accès réseau peut reproduire ce projet de A à Z en suivant les étapes décrites section par section.

\begin{tcolorbox}[
  enhanced,
  colback=PrimaryBlueLight,
  colframe=PrimaryBlue,
  boxrule=1.2pt,
  arc=8pt,
  left=12pt,right=12pt,top=10pt,bottom=10pt,
  title={\iconPerspectiveW\hspace{0.4em}\textbf{Perspective d'évolution}},
  fonttitle=\bfseries\small,
  coltitle=white,
  attach boxed title to top left={yshift=-2mm,xshift=6mm},
  boxed title style={colback=PrimaryBlue,arc=4pt,boxrule=0pt}
]
Ce projet constitue une base solide extensible : ajout de nouvelles machines (une ligne dans \texttt{hosts.yml} suffit), intégration de nouvelles applications via de nouveaux rôles Ansible, ou migration vers un environnement cloud (AWS, Azure) avec les mêmes playbooks.
\end{tcolorbox}

\vspace{3em}
\begin{center}
  \textcolor{PrimaryBlue}{\rule{4cm}{0.8pt}\hspace{1em}\textbf{\large Guardia School}\hspace{1em}\rule{4cm}{0.8pt}}\\[1em]
  \textcolor{MediumGray}{\small Équipe Projet Ansible \textperiodcentered{} Rapport finalisé le 23 avril 2026}
\end{center}

\end{document}